\documentclass[11pt]{article}
\usepackage[a4paper,margin=25mm]{geometry}
\usepackage{amsmath,amssymb,amsthm,booktabs,hyperref,listings}
\hypersetup{colorlinks=true,urlcolor=blue,linkcolor=blue}
\newtheorem{theorem}{Theorem}
\newtheorem{lemma}{Lemma}
\newcommand{\E}{\mathbb E}
\newcommand{\N}{\mathbb N}
\title{A certified numerical improvement of the BBMST minimum-modulus bound}
\author{Computational investigation and proof certificate}
\date{14 October 2026}
\begin{document}
\maketitle

\begin{abstract}
We give a computer-assisted proof that a finite covering system with distinct
positive integer moduli has a modulus at most $615299$. Equivalently, distinct
moduli all at least $615300$ cannot cover the integers. The argument retains
BBMST's distortion method, restricted second-moment estimate, and analytic
termination theorem. The changes are a modest optimization of the distortion
parameters, an efficient finite evaluation of the restricted moments, and an
explicit continuation to the analytic stopping criterion. A standalone Python
verifier uses outward-rounded integer intervals throughout. This is a local
computational proof with an explicit mathematical reduction, not a claim of
independent peer review or publication priority.
\end{abstract}

\section{Statement and scope}
The source is Balister, Bollob\'as, Morris, Sahasrabudhe and Tiba (BBMST),
\emph{On the Erd\H{o}s Covering Problem: the density of the uncovered set},
\href{https://arxiv.org/html/1811.03547v1}{arXiv:1811.03547v1}.
Numbering here refers to that version, especially Sections 6--8.
Its Theorem 8.1 excludes distinct moduli all at least $616000$.

\begin{theorem}\label{thm:main}
Every finite covering system of the integers with distinct positive integer
moduli contains a modulus at most $615299$.
\end{theorem}

This improves the exclusion cutoff by $700$. The familiar formulation ``upper
bound $616000$'' is thus replaced by the explicit integer upper bound $615299$.
We do not assert that this is the optimum attainable by the method.

The proof has two parts: the mathematical reduction below, and the finite
integer calculation in \texttt{code/verify\_minimum\_modulus.py}.
The latter reads the rational parameter witness
\texttt{data/minimum-modulus-certificate.json.gz}. It ignores the witness's
floating-point claims of success and its recorded numerical checkpoints.
The authoritative numerical output is
\texttt{data/minimum-modulus-certificate-verified.json}.

\section{The distortion criterion used}
Let $p_i$ denote the $i$th prime, including primes absent from the least common
multiple of the moduli. For each stage choose $0\leq\delta_i\leq1/2$.
Use BBMST's measures $P_i$, removed sets $B_i$, and fibre proportions $\alpha_i$.
Write
\[
 \mu_i=1-\sum_{j\leq i}P_j(B_j),\qquad
 \nu(p_i^a)=(1-\delta_i)^{-1}\quad(a\geq1),
\]
extending $\nu$ multiplicatively. Define
\[
 C_i=\prod_{j\leq i}\left(1+\frac{\nu(p_j)}{p_j-1}\right),\qquad
 T_i=\prod_{j\leq i}\left(1+
             \frac{(3p_j-1)\nu(p_j)}{(p_j-1)^2}\right).
\]
BBMST's moment bounds and measure construction imply
\begin{equation}\label{eq:loss}
 P_i(B_i)\leq
 \min\left\{M_i^{(1)},\frac{M_i^{(2)}}{4\delta_i(1-\delta_i)}\right\},
 \qquad M_i^{(r)}=\E_{i-1}\alpha_i^r,
\end{equation}
when $\delta_i>0$. With $\delta_i=0$ we use the first moment.
The factor $4\delta_i$ follows pointwise from
$(x-\delta_i)_+\leq x^2/(4\delta_i)$.
Marginal preservation makes the sum of these losses a valid upper bound for
the final covered mass.

We use the following established analytic implication, BBMST Theorem 6.1:
if $k\geq10$, $\mu_k>0$, and
\begin{equation}\label{eq:stop}
 \frac{T_k}{\mu_k}\leq k(\log k+\log\log k-3)^2,
\end{equation}
then the system does not cover the integers. Here $T_k/\mu_k$ is their $f_k$
with $i_0=0$ and $\kappa=1$. The generic second-moment bound
$M_i^{(2)}\leq T_{i-1}/(p_i-1)^2$ verifies their required hypothesis (20).
This theorem handles the entire uncomputed tail, not merely primes below a
numerical limit. If the system has no prime factors beyond a verified stage,
positivity of the survivor bound already proves noncoverage.

\section{Exact finite reduction of the restricted moments}
Fix a cutoff $K$, a current prime $p=p_i$, and the preceding primes $q<p$.
All smooth numbers in this section have prime factors strictly below $p$.
Set $C=C_{i-1}$, $T=T_{i-1}$, and, on these smooth numbers, put
\[
 w(m)=\frac{\nu(m)}m,\qquad
 F(m)=\sum_{\substack{j\geq1\\mp^j\geq K}}p^{-j},\qquad
 W(m,n)=\frac{\nu(\operatorname{lcm}(m,n))}{\operatorname{lcm}(m,n)}.
\]
The restriction that the original moduli are distinct is essential: there is
at most one progression for each $mp^j$. BBMST's indicator moment estimate
then gives
\begin{equation}\label{eq:restricted}
 M_i^{(1)}\leq\sum_m F(m)w(m),\qquad
 M_i^{(2)}\leq\sum_{m,n}F(m)F(n)W(m,n).
\end{equation}
Allowing all smooth numbers and all exponents only enlarges the sums when the
original least common multiple has restricted prime powers.

For each $j\geq1$ with $p^j<K$ define
\[
 s_j=\left\lceil K/p^j\right\rceil,
 \quad z_j=p^{-j},\quad A=\frac1{p-1}.
\]
There are finitely many such $j$, and
\begin{equation}\label{eq:F}
 F(m)=A-\sum_j z_j\,\mathbf1_{m<s_j}.
\end{equation}
All later powers have already been summed by the geometric series $A$.

Define the multiplicative function $r$ on smooth integers by
\begin{equation}\label{eq:r}
 r(q^a)=\frac{\nu(q)}{q^a}
       \frac{a+1+1/(q-1)}{1+\nu(q)/(q-1)},\qquad r(1)=1.
\end{equation}
The exact row sum of the kernel is
\begin{equation}\label{eq:row}
 \sum_n W(m,n)=C r(m).
\end{equation}
Indeed, at $q^a\Vert m$, summing over the exponent $b$ of $q$ in $n$ gives
$\nu(q)q^{-a}(a+1+1/(q-1))$; at $q\nmid m$ it gives
$1+\nu(q)/(q-1)$. Also $\sum_{m,n}W(m,n)=T$, by the same Euler-factor
calculation as BBMST Lemma 3.7.

For efficient finite pair sums, define $h(m)=m/\nu(m)$ and
$J=h*\mu_{\mathrm{Mob}}$, where $\mu_{\mathrm{Mob}}$ is the M\"obius function.
This $\mu_{\mathrm{Mob}}$ is unrelated to the survivor mass $\mu_i$.
The local factors are
\begin{equation}\label{eq:J}
 J(1)=1,\qquad J(q)=\frac q{\nu(q)}-1,\qquad
 J(q^a)=\frac{q^a-q^{a-1}}{\nu(q)}\quad(a\geq2).
\end{equation}
They are nonnegative because $\nu(q)\leq2\leq q$.
Multiplicativity gives the identity
\[
 W(m,n)=w(m)w(n)h(\gcd(m,n))
       =w(m)w(n)\sum_{d\mid m,\ d\mid n}J(d).
\]
Set $N=s_1-1$ when there are restricted powers, and $N=0$ otherwise. Extend
$w,r,J$ by zero on nonsmooth integers in the following finite sums, and define
\[
 U=C\sum_jz_j\sum_{m<s_j}r(m),\qquad
 v(m)=w(m)\sum_{j:m<s_j}z_j,
\]
\[
 V=\sum_{d=1}^{N}J(d)
          \left(\sum_{\substack{1\leq m\leq N\\d\mid m}}v(m)\right)^2.
\]
Expanding \eqref{eq:F} in \eqref{eq:restricted} now yields the exact majorants
\begin{align}
 \widehat M_i^{(1)}&=CA-\sum_jz_j\sum_{m<s_j}w(m),\label{eq:m1}\\
 \widehat M_i^{(2)}&=TA^2-2AU+V.\label{eq:m2}
\end{align}
These expressions evaluate BBMST's restricted sums, rather than postulating
a stronger universal moment inequality. The convolution replaces a quadratic
pair enumeration with divisor sums of order $N\log N$. At $K=615300$ and
$p=239$, $N=2574$, and it decreases thereafter. If $p\geq K$, then $U=V=0$.
There is no discarded smooth-number or prime-power tail.

\section{Initialization and rational parameter witness}
Set $K=615300$ and $\delta_i=0$ for $i\leq51$, so $p_{51}=233$.
The exact uniform-measure union bound gives
\begin{equation}\label{eq:init}
 \mu_{51}\geq1-
 \prod_{q\leq233}\frac q{q-1}
 +\sum_{\substack{1\leq m<K\\m\text{ is }233\text{-smooth}}}\frac1m.
\end{equation}
The verifier enumerates precisely the finite set in this formula.
Its certified lower bound exceeds $0.6541240631682086$.

The witness specifies an integer $D_i$ for every $52\leq i\leq1050000$;
the proof uses exactly
\[
 \delta_i= D_i/10^{12}.
\]
The verifier checks $0<D_i\leq5\cdot10^{11}$ before processing any stage.
No numerical optimality assertion is needed for the proof.

For reproducibility of the search, the unrounded exploratory choice was
\[
 d_i=t(p_i)(1-p_i^{-1/2})\,
 \frac{1+a_i}{1+\sqrt{1+4u_{i-1}a_i(1+a_i)/\widehat M_i^{(2)}}},
 \qquad a_i=\frac{3p_i-1}{(p_i-1)^2},
\]
followed by clipping to $(0,1/2]$ and rounding to multiples of $10^{-12}$.
Here $u_{i-1}$ is the exploratory survivor bound before the step.
The multiplier $t$ is constant on the following intervals, and equals $1$
above $630000$:
\begin{center}
\begin{tabular}{rr}
\toprule
Prime interval & Multiplier \\
\midrule
$233<p\leq400$ & $0.997072977247918$\\
$400<p\leq800$ & $1.0011180232054513$\\
$800<p\leq2000$ & $1.0047360222987132$\\
$2000<p\leq10000$ & $1.005711696549407$\\
$10000<p\leq100000$ & $1.003990673495079$\\
$100000<p\leq630000$ & $1.0019513485013882$\\
\bottomrule
\end{tabular}
\end{center}
These six parameters were selected by two passes of one-dimensional golden
section searches, minimizing the exploratory $f_{51000}$ at $K=615000$.
They were then reused at $K=615300$. Metadata containing these floating-point
multipliers is informational; only the integer list $D_i$ enters verification.

\section{Directed arithmetic and final inequalities}
Every interval in the verifier is represented by integers $[l,u]$ meaning
$[l/10^{40},u/10^{40}]$. Addition and subtraction use exact integer endpoints;
multiplication and positive-denominator division round outward using integer
floor and ceiling. In particular, cancellation in \eqref{eq:m2} is enclosed,
not hidden behind a relative-error allowance.

Starting from a lower endpoint of \eqref{eq:init}, maintain a scalar lower
bound $u_i$ for $\mu_i$ by subtracting an upper endpoint of the loss in
\eqref{eq:loss}. The verifier asserts $u_i>0$ at every step. Interval products
enclose $C_i,T_i$, and an upper bound for $f_i$ is the outward-rounded ratio
$\overline T_i/u_i$. The unrestricted product bound remains valid for the
uncomputed stages regardless of the earlier use of restricted moments.

Logarithms in \eqref{eq:stop} are also certified with integer arithmetic.
For $1\leq x\leq2$, put $z=(x-1)/(x+1)$ and use
\[
 \log x=2\sum_{j=0}^{79}\frac{z^{2j+1}}{2j+1}+R,
 \qquad 0\leq R\leq\frac{2z^{161}}{161(1-z^2)}.
\]
Range reduction $x=2^e y$, $1\leq y<2$, reduces all required logarithms to
this formula. Monotonicity allows a lower endpoint of $\log k$ as the input
for a lower bound on $\log\log k$. Thus the stopping test has an upper bound
on its left side and a lower bound on its right side.

At the final stage $k=1050000$, $p_k=16313527$, the verified output implies
the following deliberately coarser inequalities:
\begin{align*}
 \mu_k&>0.00049578893768,\\
 f_k&<191101044,\\
 k(\log k+\log\log k-3)^2&>191181612.
\end{align*}
The strict margin exceeds $80568$. Therefore \eqref{eq:stop} holds, and
BBMST Theorem 6.1 proves Theorem~\ref{thm:main}. The precise endpoint values
and a SHA-256 hash of the uncompressed witness are recorded in the verification
JSON file. The rational witness, rather than the generator's platform-dependent
floating-point output, fixes the calculation being certified.

\section{Experiments, checks, and limitations}
The baseline $K=616000$, $i_0=51$ computation gives
$\mu_{51}\approx0.654258882518139$, reproducing BBMST's displayed initial
value. Its exploratory $f_{51000}\approx5590148$ is close to, but not exactly,
their reported $5589593$. The generator uses the previous survivor bound in
the parameter formula; the paper displays a same-index survivor bound there.
Thus we do not claim a bit-for-bit reconstruction of their parameter sequence.
The discrepancy does not affect the verification of the supplied witness.

The searches varied the first distorted prime, uniform parameter scaling,
power-law corrections, transition ramps, and the six block multipliers above.
The JSON search records retain the results. Uniform scalings and transition
ramps generally worsened this numerical objective. These failed searches are
not impossibility proofs for smaller cutoffs.

At $K=615300$, the block adjustment reduced the exploratory $f_{51000}$
from approximately $5804109.37$ to $5792780.39$. The unadjusted sequence
experimentally reached the analytic criterion at $k=1744533$; the adjusted
sequence first did so at $k=1041581$. The delivered certificate continues to
the round index $1050000$ to obtain a comfortable strict margin. It does not
rely on rounded entries of BBMST's Table 1.

The first moment is evaluated exactly as a smooth Euler product minus a finite
sum. The second-moment implementation is an efficient exact reformulation
of BBMST's majorant. No stronger universal first- or second-moment inequality
has been established in this investigation, and no fundamental change of the
distortion method is being claimed.

The separate test script performs $2520$ exact-rational interval arithmetic
checks and $24$ moment checks against an independently implemented recursive
$\Theta_i(s,t)$ calculation. The latter sums exponent tails as exact geometric
series using Python fractions. These checks supplement the algebraic proof;
they do not replace it. This is not a proof-assistant formalization.

\section{Reproduction}
All scripts require only Python 3's standard library. Run from the repository
root. The primary verification does not run the optimizer:
\begin{lstlisting}[basicstyle=\ttfamily\small,breaklines=true]
python code/test_minimum_modulus.py
python code/verify_minimum_modulus.py data/minimum-modulus-certificate.json.gz --output data/minimum-modulus-certificate-verified.json
\end{lstlisting}
The certificate check took about half a minute in the local environment;
runtime and memory use depend on the machine. To repeat the exploratory
optimization and regenerate a candidate witness:
\begin{lstlisting}[basicstyle=\ttfamily\small,breaklines=true]
python code/optimize_minimum_modulus.py
python code/minimum_modulus_experiment.py --K 615300 --blocks data/minimum-modulus-optimization.json --limit 18000000 --min-stop 1050000 --output data/minimum-modulus-certificate.json.gz
python code/verify_minimum_modulus.py data/minimum-modulus-certificate.json.gz --output data/minimum-modulus-certificate-verified.json
\end{lstlisting}
A regenerated candidate must pass the verifier before it is used as a proof.
The generator's upper prime limit is only a search limit: it proves nothing
about the infinite tail unless the certified stopping inequality is satisfied.
\end{document}
